\thispagestyle{fancy}
\fancyhead[R]{A.U 2025-2026}
\fancyhead[L]{ chapitre 3: Release 1 : Fondations de la plateforme }
\chapter{Release 1 : Fondations de la plateforme Sellio}
\section*{Introduction}
Ce chapitre présente la première release, qui pose les fondations de Sellio. Le premier sprint découpe l'application existante en microservices derrière une passerelle et sécurise l'accès (JWT, rôles et permissions). Le second transforme la boutique unique en plateforme multi-boutiques : inscription des commerçants, assistant de création de boutique, vérification d'identité (KYC) et console d'administration de la plateforme.

%==================================================================
\section{Backlog du Sprint 1}
\Needspace{18\baselineskip}
Le tableau~\ref{tab:backlog-sprint1} présente le \textbf{backlog du Sprint 1}, consacré à l'architecture microservices, à l'authentification et au contrôle d'accès.

\renewcommand{\arraystretch}{1.05}
\setlength{\tabcolsep}{3pt}
\begin{longtable}{|>{\raggedright\arraybackslash}p{4cm}|>{\raggedright\arraybackslash}p{9cm}|>{\centering\arraybackslash}p{1.2cm}|}
\hline
\textbf{User Story} & \textbf{Tâches} & \textbf{Estimation} \\
\hline
\endfirsthead
\hline
\textbf{User Story} & \textbf{Tâches} & \textbf{Estimation} \\
\hline
\endhead
\hline
\multicolumn{3}{r}{\textit{Suite à la page suivante}} \\
\endfoot
\hline
\endlastfoot

\textbf{En tant qu'utilisateur, je souhaite créer un compte}
& Découper l'application NaturEssence en microservices (auth, catalog, order, marketing, analytics) et une bibliothèque partagée. & 1 \\
\cline{2-3}
& Mettre en place la passerelle Spring Cloud Gateway et ses règles de routage. & 1 \\
\cline{2-3}
& Développer l'API d'inscription avec hachage BCrypt et contrôle d'unicité de l'e-mail. & 1 \\
\cline{2-3}
& Développer l'interface d'inscription. & 1 \\
\cline{2-3}
& Tester la fonctionnalité. & 1 \\
\hline

\textbf{En tant qu'utilisateur, je souhaite me connecter}
& Développer l'API de connexion et la génération du jeton JWT signé. & 1 \\
\cline{2-3}
& Développer le filtre JWT partagé par tous les microservices. & 1 \\
\cline{2-3}
& Développer l'interface de connexion et la gestion des erreurs. & 1 \\
\cline{2-3}
& Tester la fonctionnalité. & 1 \\
\hline

\textbf{En tant qu'utilisateur, je souhaite rester connecté}
& Développer le jeton de rafraîchissement persistant et l'API \texttt{/refresh}. & 1 \\
\cline{2-3}
& Renouveler automatiquement le jeton côté client (intercepteur HTTP). & 1 \\
\hline

\textbf{En tant qu'utilisateur, je souhaite me déconnecter}
& Développer l'API de déconnexion (révocation du jeton de rafraîchissement) et le bouton associé. & 1 \\
\hline

\textbf{En tant que client, je souhaite consulter et modifier mon profil}
& Développer l'API et l'interface « Mon profil ». & 1 \\
\cline{2-3}
& Tester la fonctionnalité. & 1 \\
\hline

\textbf{En tant que commerçant, je souhaite créer des rôles et leur associer des permissions}
& Développer l'API CRUD des rôles et l'association des permissions. & 1 \\
\cline{2-3}
& Développer l'interface « Rôles et permissions » (matrice des permissions). & 1 \\
\cline{2-3}
& Tester la fonctionnalité. & 1 \\
\hline

\textbf{En tant que commerçant, je souhaite gérer les comptes de mon équipe et de mes clients}
& Développer l'API de gestion des utilisateurs (création, recherche, filtres, statut, statistiques). & 1 \\
\cline{2-3}
& Développer les interfaces « Clients », « Détail client » et « Ajouter un compte ». & 1 \\
\cline{2-3}
& Tester la fonctionnalité. & 1 \\
\hline

\textbf{En tant que système, je veux restreindre chaque endpoint selon les permissions}
& Configurer Spring Security dans chaque microservice (endpoints publics, profil, administration). & 1 \\
\cline{2-3}
& Protéger les routes du back-office côté interface (garde d'authentification). & 1 \\
\hline

\multicolumn{2}{|r|}{\textbf{Total}} & \textbf{22} \\
\hline
\caption{Backlog du Sprint 1}
\label{tab:backlog-sprint1}
\end{longtable}

\section{Diagramme de cas d'utilisation du Sprint 1}
\Needspace{20\baselineskip}
La figure~\ref{fig:cu-sprint1} présente le diagramme de cas d'utilisation du premier sprint.
\begin{figure}[H]
    \centering
    \img{0.8}{cu_sprint1.png}
    \caption{Diagramme de cas d'utilisation du Sprint 1}
    \label{fig:cu-sprint1}
\end{figure}

\section{Services web du Sprint 1}
Les services web permettent aux applications de communiquer via des API REST. Toutes les requêtes passent par la passerelle (\texttt{http://localhost:8080}), qui les route vers le service compétent. Le tableau~\ref{tab:ws-sprint1} présente les services web du Sprint 1.

\renewcommand{\arraystretch}{1.15}
\setlength{\tabcolsep}{4pt}
\begin{longtable}{|p{1.5cm}|p{6.8cm}|p{6.2cm}|}
\hline
\textbf{Méthode} & \textbf{URL} & \textbf{Description} \\
\hline
\endfirsthead
\hline
\textbf{Méthode} & \textbf{URL} & \textbf{Description} \\
\hline
\endhead
\multicolumn{3}{|c|}{\textbf{Authentification --- /api/v1/auth (auth-service)}} \\
\hline
POST & /api/v1/auth/register & Créer un compte client et retourner les jetons. \\
\hline
POST & /api/v1/auth/login & Authentifier un utilisateur et retourner le jeton d'accès et le jeton de rafraîchissement. \\
\hline
POST & /api/v1/auth/refresh & Obtenir un nouveau jeton d'accès. \\
\hline
POST & /api/v1/auth/logout & Révoquer le jeton de rafraîchissement. \\
\hline
GET & /api/v1/auth/me & Retourner l'utilisateur connecté, son rôle et sa boutique. \\
\hline
\multicolumn{3}{|c|}{\textbf{Profil --- /api/v1/profile}} \\
\hline
GET & /api/v1/profile & Consulter son profil. \\
\hline
PUT & /api/v1/profile & Modifier son profil. \\
\hline
\multicolumn{3}{|c|}{\textbf{Rôles et permissions --- /api/v1/admin/roles}} \\
\hline
GET & /api/v1/admin/roles & Lister les rôles. \\
\hline
POST & /api/v1/admin/roles & Créer un rôle. \\
\hline
PUT & /api/v1/admin/roles/\{id\} & Modifier un rôle. \\
\hline
DELETE & /api/v1/admin/roles/\{id\} & Supprimer un rôle. \\
\hline
PUT & /api/v1/admin/roles/\{id\}/permissions & Associer des permissions à un rôle. \\
\hline
GET & /api/v1/admin/roles/permissions-matrix & Obtenir la matrice rôles / permissions. \\
\hline
\multicolumn{3}{|c|}{\textbf{Utilisateurs --- /api/v1/admin/users}} \\
\hline
GET & /api/v1/admin/users & Lister les utilisateurs de la boutique. \\
\hline
POST & /api/v1/admin/users & Créer un compte. \\
\hline
PUT & /api/v1/admin/users/\{id\} & Modifier un compte. \\
\hline
PATCH & /api/v1/admin/users/\{id\}/status & Activer ou désactiver un compte. \\
\hline
GET & /api/v1/admin/users/search & Rechercher un utilisateur. \\
\hline
GET & /api/v1/admin/users/stats & Statistiques des comptes. \\
\hline
\caption{Services web du Sprint 1}
\label{tab:ws-sprint1}
\end{longtable}

\section{Les cas d'utilisation du Sprint 1}
\subsection{Cas d'utilisation « S'authentifier »}
\subsubsection{Description textuelle}
Le tableau~\ref{tab:uc-login} présente la description textuelle du cas d'utilisation « S'authentifier ».
\renewcommand{\arraystretch}{1.5}
\begin{longtable}{|p{4cm}|p{11cm}|}
\hline
\textbf{Acteurs} & Commerçant, Collaborateur, Client, Administrateur Sellio \\
\hline
\textbf{Objectif} & Accéder à la plateforme de façon sécurisée et obtenir un jeton reflétant son rôle, ses permissions et sa boutique. \\
\hline
\textbf{Pré-condition} & L'utilisateur possède un compte. \\
\hline
\textbf{Scénario principal} &
1. L'utilisateur saisit son e-mail et son mot de passe.\newline
2. L'interface envoie la requête à la passerelle, qui la route vers le service d'identité.\newline
3. Le système recherche le compte et compare le mot de passe à son empreinte BCrypt.\newline
4. Le système vérifie que le compte n'est pas bloqué.\newline
5. Le système génère un jeton d'accès JWT (24 heures) contenant l'identité, le rôle et la boutique, ainsi qu'un jeton de rafraîchissement (7 jours).\newline
6. Le système retourne les jetons et le profil de l'utilisateur.\newline
7. L'interface stocke les jetons et redirige l'utilisateur selon son rôle : console Sellio pour l'administrateur, back-office pour le commerçant, vitrine pour le client. \\
\hline
\textbf{Scénario alternatif} &
\textbf{Identifiants incorrects :} à l'étape 3, le système renvoie le message générique « Email ou mot de passe incorrect », sans préciser le champ erroné, pour ne pas faciliter l'énumération des comptes.\newline
\textbf{Compte bloqué :} à l'étape 4, le système répond 401 avec « Votre compte a été bloqué. Contactez l'équipe Sellio. »\newline
\textbf{Jeton expiré :} lors d'un appel ultérieur, l'intercepteur utilise le jeton de rafraîchissement pour obtenir un nouveau jeton sans redemander le mot de passe. \\
\hline
\caption{Description textuelle du cas d'utilisation « S'authentifier »}
\label{tab:uc-login}
\end{longtable}

\subsubsection{Diagramme de séquence système}
\Needspace{18\baselineskip}
La figure~\ref{fig:dss-login} présente le diagramme de séquence système de ce cas d'utilisation.
\begin{figure}[H]
    \centering
    \img{0.6}{dss_login.png}
    \caption{Diagramme de séquence système du cas « S'authentifier »}
    \label{fig:dss-login}
\end{figure}

\subsubsection{Diagramme de séquence objet}
\Needspace{20\baselineskip}
La figure~\ref{fig:dso-login} présente le diagramme de séquence objet : la requête traverse la passerelle, le contrôleur \texttt{AuthController}, le service \texttt{AuthService}, le gestionnaire d'authentification Spring Security, puis le dépôt \texttt{UserRepository}.
\begin{figure}[H]
    \centering
    \img{0.9}{dso_login.png}
    \caption{Diagramme de séquence objet du cas « S'authentifier »}
    \label{fig:dso-login}
\end{figure}

\section{Réalisation du Sprint 1}
\Needspace{20\baselineskip}
\textbf{Interface de connexion :} la figure~\ref{fig:ui-login} présente la page de connexion de Sellio, affichée par défaut en thème sombre avec une bascule vers le thème clair.
\begin{figure}[H]
    \centering
    \img{0.8}{ui_login.png}
    \caption{Interface de connexion à Sellio}
    \label{fig:ui-login}
\end{figure}

\Needspace{20\baselineskip}
\textbf{Interface de gestion des rôles et permissions :} la figure~\ref{fig:ui-roles} présente la matrice des permissions : le commerçant crée un rôle (par exemple « Préparateur de commandes ») et coche les modules auxquels il donne accès.
\begin{figure}[H]
    \centering
    \img{0.8}{ui_roles.png}
    \caption{Gestion des rôles et permissions}
    \label{fig:ui-roles}
\end{figure}

\Needspace{20\baselineskip}
\textbf{Interface de gestion des clients :} la figure~\ref{fig:ui-clients} présente la liste des clients de la boutique, avec recherche, filtres par statut et par segment, et accès à la fiche détaillée.
\begin{figure}[H]
    \centering
    \img{0.8}{ui_clients.png}
    \caption{Liste des clients de la boutique}
    \label{fig:ui-clients}
\end{figure}

%==================================================================
\section{Backlog du Sprint 2}
\Needspace{18\baselineskip}
Le tableau~\ref{tab:backlog-sprint2} présente le \textbf{backlog du Sprint 2}, consacré à la transformation en plateforme multi-boutiques.

\renewcommand{\arraystretch}{1.05}
\setlength{\tabcolsep}{3pt}
\begin{longtable}{|>{\raggedright\arraybackslash}p{4cm}|>{\raggedright\arraybackslash}p{9cm}|>{\centering\arraybackslash}p{1.2cm}|}
\hline
\textbf{User Story} & \textbf{Tâches} & \textbf{Estimation} \\
\hline
\endfirsthead
\hline
\textbf{User Story} & \textbf{Tâches} & \textbf{Estimation} \\
\hline
\endhead
\hline
\multicolumn{3}{r}{\textit{Suite à la page suivante}} \\
\endfoot
\hline
\endlastfoot

\textbf{En tant que visiteur, je souhaite découvrir Sellio sur une page d'accueil}
& Concevoir et développer la page d'accueil Sellio (thème sombre premium). & 1 \\
\cline{2-3}
& Développer la bascule thème sombre / thème clair. & 1 \\
\hline

\textbf{En tant que commerçant, je souhaite m'inscrire sur Sellio}
& Développer l'API d'inscription commerçant (\texttt{/register-merchant}). & 1 \\
\cline{2-3}
& Développer les interfaces d'inscription et de connexion Sellio. & 1 \\
\hline

\textbf{En tant que commerçant, je souhaite créer ma boutique avec un assistant}
& Ajouter l'entité \texttt{Shop} et rattacher utilisateurs, produits, commandes et événements à une boutique. & 1 \\
\cline{2-3}
& Filtrer toutes les requêtes par la boutique de l'utilisateur connecté (isolation multi-boutiques). & 1 \\
\cline{2-3}
& Développer l'API de création de boutique (nom, \textit{slug}, secteur, logo, couleurs). & 1 \\
\cline{2-3}
& Développer l'assistant en cinq étapes avec retour en arrière. & 1 \\
\cline{2-3}
& Tester la fonctionnalité. & 1 \\
\hline

\textbf{En tant que commerçant, je souhaite ouvrir ma vitrine depuis le back-office}
& Développer le lien « Voir le site » vers la vitrine de la boutique. & 1 \\
\hline

\textbf{En tant que commerçant, je souhaite enregistrer ma carte bancaire}
& Développer l'API Stripe de création d'un \textit{SetupIntent}. & 1 \\
\cline{2-3}
& Intégrer Stripe Elements et confirmer la carte côté navigateur. & 1 \\
\cline{2-3}
& Stocker uniquement la marque, les 4 derniers chiffres, l'expiration et la référence \texttt{pm\_}. & 1 \\
\hline

\textbf{En tant que commerçant, je souhaite soumettre ma pièce d'identité et un selfie}
& Développer l'entité \texttt{MerchantVerification} et l'API de soumission avec contrôles (CIN à 8 chiffres, passeport, taille des images). & 1 \\
\cline{2-3}
& Développer l'étape « Identité » (CIN recto-verso ou passeport, selfie). & 1 \\
\cline{2-3}
& Développer la page de suivi du dossier et la nouvelle soumission après refus. & 1 \\
\hline

\textbf{En tant qu'administrateur Sellio, je souhaite valider ou refuser un dossier}
& Développer les API d'approbation et de refus (motif obligatoire). & 1 \\
\cline{2-3}
& Développer l'écran de revue du dossier dans la console. & 1 \\
\cline{2-3}
& Tester la fonctionnalité. & 1 \\
\hline

\textbf{En tant qu'administrateur Sellio, je souhaite consulter les boutiques et leurs propriétaires}
& Créer le compte administrateur Sellio au démarrage (\textit{seed}). & 1 \\
\cline{2-3}
& Développer les API de statistiques, de liste et de détail des boutiques. & 1 \\
\cline{2-3}
& Développer la console avec pagination. & 1 \\
\hline

\textbf{En tant qu'administrateur Sellio, je souhaite suspendre une boutique ou bloquer un commerçant}
& Développer les API de changement de statut d'une boutique et d'un utilisateur. & 1 \\
\cline{2-3}
& Refuser la connexion d'un compte bloqué et afficher un message de boutique fermée sur la vitrine. & 1 \\
\cline{2-3}
& Tester la fonctionnalité. & 1 \\
\hline

\multicolumn{2}{|r|}{\textbf{Total}} & \textbf{25} \\
\hline
\caption{Backlog du Sprint 2}
\label{tab:backlog-sprint2}
\end{longtable}

\section{Diagramme de cas d'utilisation du Sprint 2}
\Needspace{20\baselineskip}
La figure~\ref{fig:cu-sprint2} présente le diagramme de cas d'utilisation du deuxième sprint.
\begin{figure}[H]
    \centering
    \img{0.85}{cu_sprint2.png}
    \caption{Diagramme de cas d'utilisation du Sprint 2}
    \label{fig:cu-sprint2}
\end{figure}

\section{Services web du Sprint 2}
Le tableau~\ref{tab:ws-sprint2} présente les services web du Sprint 2.
\renewcommand{\arraystretch}{1.15}
\setlength{\tabcolsep}{4pt}
\begin{longtable}{|p{1.5cm}|p{6.8cm}|p{6.2cm}|}
\hline
\textbf{Méthode} & \textbf{URL} & \textbf{Description} \\
\hline
\endfirsthead
\hline
\textbf{Méthode} & \textbf{URL} & \textbf{Description} \\
\hline
\endhead
\multicolumn{3}{|c|}{\textbf{Commerçant et boutique --- /api/v1/auth}} \\
\hline
POST & /api/v1/auth/register-merchant & Créer un compte commerçant. \\
\hline
POST & /api/v1/auth/my-shop & Créer sa boutique avec son dossier de vérification. \\
\hline
GET & /api/v1/auth/my-shop & Consulter sa boutique. \\
\hline
PATCH & /api/v1/auth/my-shop & Modifier les informations de sa boutique. \\
\hline
GET & /api/v1/auth/my-shop/verification & Consulter l'état de son dossier. \\
\hline
PUT & /api/v1/auth/my-shop/verification & Soumettre à nouveau un dossier refusé. \\
\hline
GET & /api/v1/public/shops/\{slug\} & Informations publiques d'une boutique (statut, gabarit, thème). \\
\hline
POST & /api/v1/public/stripe/setup-intent & Créer un \textit{SetupIntent} Stripe pour enregistrer la carte. \\
\hline
\multicolumn{3}{|c|}{\textbf{Console plateforme --- /api/v1/admin/platform}} \\
\hline
GET & /api/v1/admin/platform/stats & Statistiques globales (boutiques, commerçants, dossiers en attente). \\
\hline
GET & /api/v1/admin/platform/shops & Lister les boutiques et leurs propriétaires (paginé). \\
\hline
GET & /api/v1/admin/platform/shops/\{id\} & Détail d'une boutique. \\
\hline
PATCH & /api/v1/admin/platform/shops/\{id\}/status & Suspendre ou réactiver une boutique. \\
\hline
PATCH & /api/v1/admin/platform/users/\{id\}/status & Bloquer ou débloquer un commerçant. \\
\hline
GET & /api/v1/admin/platform/users & Lister les utilisateurs de la plateforme. \\
\hline
GET & /api/v1/admin/platform/verifications & Lister les dossiers de vérification. \\
\hline
GET & /api/v1/admin/platform/verifications/\{id\} & Consulter un dossier (pièces et selfie). \\
\hline
POST & /api/v1/admin/platform/verifications/\{id\}/approve & Valider un dossier : la boutique passe à \texttt{ACTIVE}. \\
\hline
POST & /api/v1/admin/platform/verifications/\{id\}/reject & Refuser un dossier avec un motif : la boutique passe à \texttt{REJECTED}. \\
\hline
\caption{Services web du Sprint 2}
\label{tab:ws-sprint2}
\end{longtable}

\section{Les cas d'utilisation du Sprint 2}
\subsection{Cas d'utilisation « Créer une boutique et faire vérifier son identité »}
\subsubsection{Description textuelle}
Le tableau~\ref{tab:uc-kyc} présente la description textuelle de ce cas d'utilisation.
\renewcommand{\arraystretch}{1.5}
\begin{longtable}{|p{4cm}|p{11cm}|}
\hline
\textbf{Acteurs} & Commerçant, Administrateur Sellio, Stripe \\
\hline
\textbf{Objectif} & Ouvrir une boutique qui ne pourra vendre qu'après vérification de l'identité et du moyen de paiement de son propriétaire. \\
\hline
\textbf{Pré-condition} & Le commerçant est inscrit et connecté ; il ne possède pas encore de boutique. \\
\hline
\textbf{Scénario principal} &
1. Le commerçant saisit le nom de sa boutique ; le système propose un identifiant (\textit{slug}).\newline
2. Il choisit le secteur (cosmétiques ou vêtements), puis son logo et sa palette.\newline
3. À l'étape « Carte », le système crée un \textit{SetupIntent} Stripe ; le commerçant saisit sa carte dans le champ sécurisé Stripe, qui la confirme sans que le numéro ne transite par Sellio.\newline
4. À l'étape « Identité », il choisit CIN ou passeport, photographie sa pièce et prend un selfie.\newline
5. Le système contrôle le dossier (format du numéro, validité de la carte, taille des images), crée la boutique au statut \texttt{PENDING} et le dossier au statut en attente.\newline
6. Le back-office affiche « Dossier en cours de vérification » ; la vitrine n'est pas encore ouverte au public.\newline
7. L'administrateur Sellio consulte le dossier, compare le selfie à la pièce d'identité et le valide.\newline
8. La boutique passe au statut \texttt{ACTIVE} ; le commerçant accède à tout le back-office. \\
\hline
\textbf{Scénario alternatif} &
\textbf{Carte refusée par Stripe :} à l'étape 3, le message de Stripe est affiché et le commerçant peut saisir une autre carte.\newline
\textbf{Dossier incomplet :} à l'étape 5, le système indique la pièce manquante (par exemple « Le numéro de CIN doit contenir 8 chiffres »).\newline
\textbf{Dossier refusé :} à l'étape 7, l'administrateur saisit obligatoirement un motif ; la boutique passe à \texttt{REJECTED}, le motif est affiché au commerçant qui peut soumettre un nouveau dossier. \\
\hline
\caption{Description textuelle du cas « Créer une boutique et faire vérifier son identité »}
\label{tab:uc-kyc}
\end{longtable}

\subsubsection{Diagramme de séquence système}
\Needspace{18\baselineskip}
La figure~\ref{fig:dss-kyc} présente le diagramme de séquence système de ce cas.
\begin{figure}[H]
    \centering
    \img{0.8}{dss_kyc.png}
    \caption{Diagramme de séquence système du cas « Créer une boutique et faire vérifier son identité »}
    \label{fig:dss-kyc}
\end{figure}

\subsubsection{Diagramme de séquence objet}
\Needspace{20\baselineskip}
La figure~\ref{fig:dso-kyc} présente le diagramme de séquence objet de ce cas.
\begin{figure}[H]
    \centering
    \img{0.95}{dso_kyc.png}
    \caption{Diagramme de séquence objet du cas « Créer une boutique et faire vérifier son identité »}
    \label{fig:dso-kyc}
\end{figure}

\section{Réalisation du Sprint 2}
\Needspace{20\baselineskip}
\textbf{Page d'accueil Sellio :} la figure~\ref{fig:ui-sellio-home} présente la page d'accueil de la plateforme, qui présente l'offre aux commerçants et donne accès à l'inscription et à la connexion.
\begin{figure}[H]
    \centering
    \img{0.85}{ui_sellio_home.png}
    \caption{Page d'accueil de Sellio}
    \label{fig:ui-sellio-home}
\end{figure}

\Needspace{20\baselineskip}
\textbf{Assistant de création de boutique :} les figures~\ref{fig:ui-wizard-1} et~\ref{fig:ui-wizard-2} présentent l'assistant : identité de la boutique et secteur, puis enregistrement de la carte et de la pièce d'identité.
\begin{figure}[H]
    \centering
    \img{0.8}{ui_nouvelle_boutique.png}
    \caption{Assistant de création de boutique : informations et secteur}
    \label{fig:ui-wizard-1}
\end{figure}
\begin{figure}[H]
    \centering
    \img{0.8}{ui_kyc_carte_identite.png}
    \caption{Assistant de création de boutique : carte bancaire et pièce d'identité}
    \label{fig:ui-wizard-2}
\end{figure}

\Needspace{20\baselineskip}
\textbf{Console d'administration Sellio :} la figure~\ref{fig:ui-console} présente la console : indicateurs globaux, liste paginée des boutiques avec leur propriétaire et leur statut, et actions de suspension ou de blocage. La figure~\ref{fig:ui-verif} présente la revue d'un dossier de vérification.
\begin{figure}[H]
    \centering
    \img{0.85}{ui_console_sellio.png}
    \caption{Console d'administration Sellio}
    \label{fig:ui-console}
\end{figure}
\begin{figure}[H]
    \centering
    \img{0.85}{ui_revue_dossier.png}
    \caption{Revue d'un dossier de vérification par l'administrateur}
    \label{fig:ui-verif}
\end{figure}

\subsection*{Conclusion}
Au terme de cette release, Sellio dispose d'une architecture microservices sécurisée et d'un véritable modèle multi-boutiques : un commerçant peut s'inscrire, créer sa boutique et la faire vérifier, et l'administrateur de la plateforme garde le contrôle sur les boutiques ouvertes. La release suivante porte sur le catalogue et la vitrine.
